\section{Теоретическая часть}

\subsection{Нелинейные элементы и аппроксимация ВАХ}

Нелинейной называют электрическую цепь, в которой не выполняется принцип суперпозиции. Свойства нелинейного элемента описывают вольт-амперной характеристикой (ВАХ) $i=f(u)$. Для анализа её заменяют приближённой зависимостью. При небольших отклонениях напряжения от рабочей точки $U_0$ используют полиномиальную аппроксимацию:
\[
i(U_0+u)=i(U_0)+a_1u+a_2u^2+\ldots+a_nu^n,
\qquad a_k=\frac{1}{k!}\left.\frac{d^ki}{dU^k}\right|_{U=U_0}.
\]
Здесь $u$ --- переменная составляющая напряжения, $a_1$ --- крутизна ВАХ в рабочей точке. Для больших сигналов применяют кусочно-линейную аппроксимацию, при которой характеристику заменяют отрезками прямых. Простейшая модель диода имеет вид
\[
i(u)=\begin{cases}0,&u\leq U_1,\\S(u-U_1),&u>U_1,\end{cases}
\]
где $U_1$ --- напряжение начала проводимости, $S$ --- крутизна открытого участка.

Если на нелинейный элемент подать гармоническое напряжение $u=E\cos(\omega t+\varphi)$, после подстановки в полином и преобразования степеней косинуса получим
\[
i(t)=I_0+\sum_{k=1}^{n}I_k\cos\bigl(k\omega t+k\varphi\bigr).
\]
Таким образом, помимо основной частоты возникают постоянная составляющая и высшие гармоники. Чётные степени полинома определяют постоянную составляющую и чётные гармоники, нечётные --- нечётные гармоники.

\subsection{Комбинационные частоты и выделение полезного сигнала}

При двухчастотном воздействии
\[
u(t)=E_1\cos\omega_1t+E_2\cos\omega_2t
\]
нелинейный элемент создаёт не только гармоники исходных частот, но и комбинационные составляющие. Например, квадратичный член ВАХ даёт
\begin{align*}
a_2u^2(t)={}&\frac{a_2}{2}(E_1^2+E_2^2)
+\frac{a_2E_1^2}{2}\cos2\omega_1t
+\frac{a_2E_2^2}{2}\cos2\omega_2t\\
&+a_2E_1E_2\bigl[\cos(\omega_1+\omega_2)t+\cos(\omega_1-\omega_2)t\bigr].
\end{align*}
Следовательно, в спектре появляются суммарная $\omega_1+\omega_2$ и разностная $|\omega_1-\omega_2|$ частоты. При учёте членов более высокого порядка возможны комбинации $|m\omega_1\pm n\omega_2|$. Именно это свойство используется при преобразовании частоты.

Для выделения нужной составляющей к нелинейному элементу подключают частотно-избирательную нагрузку. Параллельная $RC$-цепь выделяет постоянную и низкочастотные составляющие, поскольку её импеданс
\[
|Z_{RC}|=\frac{R}{\sqrt{1+(\omega RC)^2}}
\]
уменьшается с ростом частоты. Параллельный $LC$-контур, напротив, выделяет колебания вблизи резонансной частоты
\[
\omega_0=\frac{1}{\sqrt{LC}}.
\]
На резонансе сопротивление контура и создаваемое на нём напряжение максимальны. Ширина полосы выделяемых частот определяется добротностью контура.

\subsection{Преобразование частоты}

Для преобразования частоты на нелинейный элемент одновременно подают исследуемый сигнал частоты $f_c$ и колебание гетеродина частоты $f_{\text{г}}$. В результате возникают комбинационные составляющие. При преобразовании по первой гармонике гетеродина разностная частота равна
\[
f_{\text{пч}}=|f_c-f_{\text{г}}|.
\]
Если нагрузкой служит $RC$-цепь, высокочастотные компоненты подавляются, а разностная составляющая наблюдается как биения. При $f_c=f_{\text{г}}$ частота биений обращается в нуль.

При резонансной нагрузке выделяются только те комбинационные составляющие, частота которых близка к частоте настройки контура $f_{\text{рез}}$. Для первой гармоники гетеродина максимумы соответствуют
\[
f_c=f_{\text{г}}\pm f_{\text{рез}}.
\]
Кроме них возможен прямой канал $f_c=f_{\text{рез}}$ и побочные каналы, связанные с высшими гармониками гетеродина, например $f_c=2f_{\text{г}}\pm f_{\text{рез}}$. Поэтому частотная характеристика преобразователя может содержать несколько максимумов, хотя полезным является лишь один канал.

\subsection{Амплитудное детектирование}

Амплитудное детектирование --- выделение низкочастотного модулирующего сигнала из высокочастотного колебания. Амплитудно-модулированное напряжение имеет вид
\[
u(t)=E_0\bigl[1+M\cos(\Omega t)\bigr]\cos(\omega_0t),
\qquad \Omega\ll\omega_0,
\]
где $E_0$ --- амплитуда несущей, $M$ --- глубина модуляции, $\Omega$ --- круговая частота модуляции, $\omega_0$ --- круговая частота несущей. В спектре такого сигнала присутствуют несущая и две боковые составляющие с частотами $\omega_0\pm\Omega$.

В диодном детекторе ток проходит преимущественно около положительных максимумов несущей, заряжая конденсатор. Между импульсами тока конденсатор разряжается через сопротивление нагрузки, вследствие чего выходное напряжение повторяет огибающую входного колебания. Для сглаживания несущей без заметного искажения огибающей постоянная времени должна удовлетворять условию
\[
\frac{2\pi}{\omega_0}\ll RC\ll\frac{2\pi}{\Omega}.
\]
При слишком малом $RC$ остаются высокочастотные пульсации, при слишком большом конденсатор не успевает разряжаться вслед за огибающей, что приводит к искажениям.

Детекторная характеристика представляет собой зависимость постоянной составляющей выходного напряжения от амплитуды входного сигнала: $U_{\text{вых}}=\varphi(U_{\text{вх}})$. При малых амплитудах она близка к квадратичной ($U_{\text{вых}}\sim U_{\text{вх}}^2$), а при больших --- к линейной. Частотная характеристика детектора показывает зависимость амплитуды выделенного низкочастотного сигнала от частоты модуляции. Изменение ёмкости нагрузки меняет постоянную времени $RC$ и, следовательно, степень сглаживания и искажения огибающей.

\subsection{Влияние детектора на резонансную цепь}

Подключённый к колебательному контуру диодный детектор представляет собой дополнительную нагрузку. Его входное сопротивление шунтирует контур, увеличивает потери и обычно уменьшает резонансное напряжение и добротность. Поэтому при подключённом детекторе резонансная кривая может становиться ниже и шире. Сравнение характеристик контура при подключённом и отключённом детекторе позволяет оценить влияние нагрузки на его избирательность.
